\chapter{Mini-Examenes de Vectores}

\section{[C06.S01] Bases y dependencia en \texorpdfstring{\(\R^2\)}{R2}}

\begin{exambrief}{Mini-examen C06.S01}
Duracion sugerida: \(25\) minutos. Valor orientativo: \(10\) puntos.
\begin{enumerate}
  \item Estudia si \(\vec u=(2,1)\) y \(\vec v=(1,3)\) forman una base de \(\R^2\).
  \item Si la respuesta es afirmativa, expresa \(\vec d=(13,14)\) como combinacion lineal de
  \(\vec u\) y \(\vec v\).
\end{enumerate}
\end{exambrief}

\begin{shortanswer}{Mini-examen C06.S01}
Si, porque
\[
\det\begin{pmatrix}2&1\\1&3\end{pmatrix}=5\neq 0.
\]
Ademas,
\[
\vec d=5\vec u+3\vec v.
\]
\end{shortanswer}

\begin{fullsolution}{Mini-examen C06.S01}
Los vectores forman una base si son linealmente independientes. Calculamos
\[
\det\begin{pmatrix}2&1\\1&3\end{pmatrix}=6-1=5\neq 0.
\]
Luego forman una base de \(\R^2\).

Buscamos \(a,b\) tales que
\[
a(2,1)+b(1,3)=(13,14).
\]
Eso da el sistema
\[
\begin{cases}
2a+b=13,\\
a+3b=14.
\end{cases}
\]
De la primera ecuacion, \(b=13-2a\). Sustituyendo:
\[
a+3(13-2a)=14
\Longrightarrow
a+39-6a=14
\Longrightarrow
-5a=-25
\Longrightarrow
a=5.
\]
Entonces
\[
b=13-10=3.
\]
\end{fullsolution}

\section{[C06.S02] Operaciones con vectores}

\begin{exambrief}{Mini-examen C06.S02}
Duracion sugerida: \(25\) minutos. Valor orientativo: \(10\) puntos.
\begin{enumerate}
  \item Dados \(\vec u=(4,-2)\) y \(\vec v=(-1,5)\), calcula \(\vec u+\vec v\).
  \item Calcula \(\vec u-\vec v\) e interpreta cual de los dos resultados representa una
  correccion respecto del otro.
\end{enumerate}
\end{exambrief}

\begin{shortanswer}{Mini-examen C06.S02}
\[
\vec u+\vec v=(3,3),
\qquad
\vec u-\vec v=(5,-7).
\]
\end{shortanswer}

\begin{fullsolution}{Mini-examen C06.S02}
Sumando coordenadas:
\[
\vec u+\vec v=(4,-2)+(-1,5)=(3,3).
\]
Restando:
\[
\vec u-\vec v=(4,-2)-(-1,5)=(5,-7).
\]
El vector diferencia representa la correccion necesaria para pasar del desplazamiento \(\vec v\)
al desplazamiento \(\vec u\).
\end{fullsolution}

\section{[C06.S03] Modulo, argumento y vectores unitarios}

\begin{exambrief}{Mini-examen C06.S03}
Duracion sugerida: \(25\) minutos. Valor orientativo: \(10\) puntos.
\begin{enumerate}
  \item Halla el modulo y el argumento del vector \(\vec v=(-8,8\sqrt3)\).
  \item Obtén el vector unitario asociado.
\end{enumerate}
\end{exambrief}

\begin{shortanswer}{Mini-examen C06.S03}
\[
\|\vec v\|=16,
\qquad
\arg(\vec v)=120^\circ=\frac{2\pi}{3},
\]
y el unitario es
\[
\left(-\frac12,\frac{\sqrt3}{2}\right).
\]
\end{shortanswer}

\begin{fullsolution}{Mini-examen C06.S03}
\[
\|\vec v\|=\sqrt{(-8)^2+(8\sqrt3)^2}=\sqrt{64+192}=\sqrt{256}=16.
\]
El vector esta en el segundo cuadrante y su angulo de referencia satisface
\[
\tg\beta=\sqrt3,
\]
de modo que \(\beta=60^\circ\). Luego
\[
\arg(\vec v)=180^\circ-60^\circ=120^\circ=\frac{2\pi}{3}.
\]
El vector unitario es
\[
\frac{\vec v}{16}=\left(-\frac12,\frac{\sqrt3}{2}\right).
\]
\end{fullsolution}

\section{[C06.S04] Producto escalar y angulo entre vectores}

\begin{exambrief}{Mini-examen C06.S04}
Duracion sugerida: \(30\) minutos. Valor orientativo: \(10\) puntos.
\begin{enumerate}
  \item Calcula el producto escalar de \(\vec u=(4,1)\) y \(\vec v=(2,5)\).
  \item Halla el angulo entre ambos vectores.
\end{enumerate}
\end{exambrief}

\begin{shortanswer}{Mini-examen C06.S04}
\[
\vec u\cdot \vec v=13,
\qquad
\cos\theta=\frac{13}{\sqrt{17}\sqrt{29}}.
\]
El angulo es aproximadamente \(54.2^\circ\).
\end{shortanswer}

\begin{fullsolution}{Mini-examen C06.S04}
\[
\vec u\cdot \vec v=4\cdot 2+1\cdot 5=13.
\]
Los modulos son
\[
\|\vec u\|=\sqrt{17},
\qquad
\|\vec v\|=\sqrt{29}.
\]
Por la formula del producto escalar:
\[
\cos\theta=\frac{13}{\sqrt{17}\sqrt{29}}.
\]
Por tanto,
\[
\theta\approx \arccos\left(\frac{13}{\sqrt{17}\sqrt{29}}\right)\approx 54.2^\circ.
\]
\end{fullsolution}

\section{[C06.S05] Paralelismo, perpendicularidad y alineacion}

\begin{exambrief}{Mini-examen C06.S05}
Duracion sugerida: \(25\) minutos. Valor orientativo: \(10\) puntos.
\begin{enumerate}
  \item Comprueba si \(A(2,1)\), \(B(6,4)\) y \(C(10,7)\) estan alineados.
  \item Estudia si el segmento \(CD\), con \(D(13,3)\), es perpendicular a la recta anterior.
\end{enumerate}
\end{exambrief}

\begin{shortanswer}{Mini-examen C06.S05}
Si, porque
\[
\overrightarrow{AB}=(4,3),\qquad \overrightarrow{BC}=(4,3).
\]
Ademas,
\[
\overrightarrow{CD}=(3,-4),
\]
y
\[
(4,3)\cdot(3,-4)=0.
\]
\end{shortanswer}

\begin{fullsolution}{Mini-examen C06.S05}
\[
\overrightarrow{AB}=(6-2,4-1)=(4,3),
\qquad
\overrightarrow{BC}=(10-6,7-4)=(4,3).
\]
Como ambos vectores son iguales, los puntos estan alineados.

Por otra parte,
\[
\overrightarrow{CD}=(13-10,3-7)=(3,-4).
\]
El producto escalar es
\[
(4,3)\cdot(3,-4)=12-12=0.
\]
Luego el segmento \(CD\) es perpendicular a la recta que contiene a \(A\), \(B\) y \(C\).
\end{fullsolution}

\section{[C06.S06] Construcciones con vectores y paralelogramos}

\begin{exambrief}{Mini-examen C06.S06}
Duracion sugerida: \(30\) minutos. Valor orientativo: \(10\) puntos.
\begin{enumerate}
  \item En un paralelogramo se conocen \(A(0,0)\), \(B(5,1)\) y \(D(2,4)\), siendo \(AB\) y
  \(AD\) lados consecutivos. Halla \(C\).
  \item Calcula el punto de corte de las diagonales y el area del paralelogramo.
\end{enumerate}
\end{exambrief}

\begin{shortanswer}{Mini-examen C06.S06}
\[
C=(7,5),
\qquad
M=\left(\frac72,\frac52\right).
\]
El area del paralelogramo es \(18\).
\end{shortanswer}

\begin{fullsolution}{Mini-examen C06.S06}
En un paralelogramo,
\[
C=B+D-A.
\]
Por tanto,
\[
C=(5,1)+(2,4)-(0,0)=(7,5).
\]
El punto medio de las diagonales es
\[
M=\frac{A+C}{2}=
\left(\frac{0+7}{2},\frac{0+5}{2}\right)=
\left(\frac72,\frac52\right).
\]
Para el area usamos
\[
\overrightarrow{AB}=(5,1),
\qquad
\overrightarrow{AD}=(2,4).
\]
Entonces
\[
Area=
\left|
\det
\begin{pmatrix}
5&1\\
2&4
\end{pmatrix}
\right|
=|20-2|=18.
\]
\end{fullsolution}
